% Options for packages loaded elsewhere
\PassOptionsToPackage{unicode}{hyperref}
\PassOptionsToPackage{hyphens}{url}
\documentclass[
  ignorenonframetext,
]{beamer}
\newif\ifbibliography
\usepackage{pgfpages}
\setbeamertemplate{caption}[numbered]
\setbeamertemplate{caption label separator}{: }
\setbeamercolor{caption name}{fg=normal text.fg}
\beamertemplatenavigationsymbolsempty
% remove section numbering
\setbeamertemplate{part page}{
  \centering
  \begin{beamercolorbox}[sep=16pt,center]{part title}
    \usebeamerfont{part title}\insertpart\par
  \end{beamercolorbox}
}
\setbeamertemplate{section page}{
  \centering
  \begin{beamercolorbox}[sep=12pt,center]{section title}
    \usebeamerfont{section title}\insertsection\par
  \end{beamercolorbox}
}
\setbeamertemplate{subsection page}{
  \centering
  \begin{beamercolorbox}[sep=8pt,center]{subsection title}
    \usebeamerfont{subsection title}\insertsubsection\par
  \end{beamercolorbox}
}
% Prevent slide breaks in the middle of a paragraph
\widowpenalties 1 10000
\raggedbottom
\AtBeginPart{
  \frame{\partpage}
}
\AtBeginSection{
  \ifbibliography
  \else
    \frame{\sectionpage}
  \fi
}
\AtBeginSubsection{
  \frame{\subsectionpage}
}
\usepackage{iftex}
\ifPDFTeX
  \usepackage[T1]{fontenc}
  \usepackage[utf8]{inputenc}
  \usepackage{textcomp} % provide euro and other symbols
\else % if luatex or xetex
  \usepackage{unicode-math} % this also loads fontspec
  \defaultfontfeatures{Scale=MatchLowercase}
  \defaultfontfeatures[\rmfamily]{Ligatures=TeX,Scale=1}
\fi
\usepackage{lmodern}
\ifPDFTeX\else
  % xetex/luatex font selection
\fi
% Use upquote if available, for straight quotes in verbatim environments
\IfFileExists{upquote.sty}{\usepackage{upquote}}{}
\IfFileExists{microtype.sty}{% use microtype if available
  \usepackage[]{microtype}
  \UseMicrotypeSet[protrusion]{basicmath} % disable protrusion for tt fonts
}{}
\makeatletter
\@ifundefined{KOMAClassName}{% if non-KOMA class
  \IfFileExists{parskip.sty}{%
    \usepackage{parskip}
  }{% else
    \setlength{\parindent}{0pt}
    \setlength{\parskip}{6pt plus 2pt minus 1pt}}
}{% if KOMA class
  \KOMAoptions{parskip=half}}
\makeatother
\usepackage{color}
\usepackage{fancyvrb}
\newcommand{\VerbBar}{|}
\newcommand{\VERB}{\Verb[commandchars=\\\{\}]}
\DefineVerbatimEnvironment{Highlighting}{Verbatim}{commandchars=\\\{\}}
% Add ',fontsize=\small' for more characters per line
\newenvironment{Shaded}{}{}
\newcommand{\AlertTok}[1]{\textcolor[rgb]{1.00,0.00,0.00}{\textbf{#1}}}
\newcommand{\AnnotationTok}[1]{\textcolor[rgb]{0.38,0.63,0.69}{\textbf{\textit{#1}}}}
\newcommand{\AttributeTok}[1]{\textcolor[rgb]{0.49,0.56,0.16}{#1}}
\newcommand{\BaseNTok}[1]{\textcolor[rgb]{0.25,0.63,0.44}{#1}}
\newcommand{\BuiltInTok}[1]{\textcolor[rgb]{0.00,0.50,0.00}{#1}}
\newcommand{\CharTok}[1]{\textcolor[rgb]{0.25,0.44,0.63}{#1}}
\newcommand{\CommentTok}[1]{\textcolor[rgb]{0.38,0.63,0.69}{\textit{#1}}}
\newcommand{\CommentVarTok}[1]{\textcolor[rgb]{0.38,0.63,0.69}{\textbf{\textit{#1}}}}
\newcommand{\ConstantTok}[1]{\textcolor[rgb]{0.53,0.00,0.00}{#1}}
\newcommand{\ControlFlowTok}[1]{\textcolor[rgb]{0.00,0.44,0.13}{\textbf{#1}}}
\newcommand{\DataTypeTok}[1]{\textcolor[rgb]{0.56,0.13,0.00}{#1}}
\newcommand{\DecValTok}[1]{\textcolor[rgb]{0.25,0.63,0.44}{#1}}
\newcommand{\DocumentationTok}[1]{\textcolor[rgb]{0.73,0.13,0.13}{\textit{#1}}}
\newcommand{\ErrorTok}[1]{\textcolor[rgb]{1.00,0.00,0.00}{\textbf{#1}}}
\newcommand{\ExtensionTok}[1]{#1}
\newcommand{\FloatTok}[1]{\textcolor[rgb]{0.25,0.63,0.44}{#1}}
\newcommand{\FunctionTok}[1]{\textcolor[rgb]{0.02,0.16,0.49}{#1}}
\newcommand{\ImportTok}[1]{\textcolor[rgb]{0.00,0.50,0.00}{\textbf{#1}}}
\newcommand{\InformationTok}[1]{\textcolor[rgb]{0.38,0.63,0.69}{\textbf{\textit{#1}}}}
\newcommand{\KeywordTok}[1]{\textcolor[rgb]{0.00,0.44,0.13}{\textbf{#1}}}
\newcommand{\NormalTok}[1]{#1}
\newcommand{\OperatorTok}[1]{\textcolor[rgb]{0.40,0.40,0.40}{#1}}
\newcommand{\OtherTok}[1]{\textcolor[rgb]{0.00,0.44,0.13}{#1}}
\newcommand{\PreprocessorTok}[1]{\textcolor[rgb]{0.74,0.48,0.00}{#1}}
\newcommand{\RegionMarkerTok}[1]{#1}
\newcommand{\SpecialCharTok}[1]{\textcolor[rgb]{0.25,0.44,0.63}{#1}}
\newcommand{\SpecialStringTok}[1]{\textcolor[rgb]{0.73,0.40,0.53}{#1}}
\newcommand{\StringTok}[1]{\textcolor[rgb]{0.25,0.44,0.63}{#1}}
\newcommand{\VariableTok}[1]{\textcolor[rgb]{0.10,0.09,0.49}{#1}}
\newcommand{\VerbatimStringTok}[1]{\textcolor[rgb]{0.25,0.44,0.63}{#1}}
\newcommand{\WarningTok}[1]{\textcolor[rgb]{0.38,0.63,0.69}{\textbf{\textit{#1}}}}
\usepackage{longtable,booktabs,array}
\usepackage{caption}
\captionsetup[table]{skip=6pt}
\newcounter{none} % for unnumbered tables
\usepackage{calc} % for calculating minipage widths
\usepackage{caption}
% Make caption package work with longtable
\makeatletter
\def\fnum@table{\tablename~\thetable}
\makeatother
\setlength{\emergencystretch}{3em} % prevent overfull lines
\providecommand{\tightlist}{%
  \setlength{\itemsep}{0pt}\setlength{\parskip}{0pt}}
% Manuscript macro/package fallbacks (newcommand -> providecommand).
\usepackage{amsmath}
\usepackage{amssymb}
\usepackage{booktabs}
% Contents listing: article.cls prefixes every \l@section entry with
% \addvspace{1.0em}, which across ten sections costs about ten lines and pushed
% the ~40-entry listing one line past page 2 — leaving a third sheet carrying a
% single "10 References" line. Tighten that inter-section gap for the duration
% of \tableofcontents only, inside a group, so body spacing is untouched and no
% entry is dropped (reducing \tocdepth would have hidden 29 real subsections).
  \providecommand{\addvspace}[1]{\vskip 2\p@}%
% Document metadata. config.yaml declares a subtitle and a keyword list, and
% the producer emits both as tokens, but the template's \hypersetup writes only
% pdftitle/pdfauthor/pdflang -- so `pdfinfo` reported no Subject and no
% Keywords. This block runs after the template's, so it wins.
%
% The two values below are WRITTEN by scripts/z_generate_manuscript_variables.py
% (sync_preamble_metadata), never typed. A double-brace token cannot be used here:
% preamble.md is substituted into output/manuscript/ like any other section, but
% the template's _manuscript_source.py then copies the raw file back over that
% copy, so an unresolved token would reach hyperref and land in the PDF's
% metadata verbatim.
% Bounded identifier splitting. The template defines
%   \protected\def\breaktt#1{\begingroup\ttfamily\seqsplit{#1}\endgroup}
% and \seqsplit permits a break after EVERY character with no continuation cue.
% In the 2026-09-05 PDF that rendered `model_family_manifest.json` as
% `model_family_manifest.js` + `on` and `src/manuscript_variables.py` as
% `(src/manusc` + `ript_variables.py` — the first is itself a valid filename, so
% a reader cannot tell a line break from a name. Keep seqsplit's guarantee that
% no code token overflows the measure, but forbid breaks inside the last
% \GNNttTail characters (an extension is never orphaned) and inside the first
% \GNNttHead (a one- or two-letter stub is never left behind).

% Slide layout defaults for warning-clean scientific decks.
\usepackage{etoolbox}
\IfFileExists{xurl.sty}{\usepackage{xurl}}{}
\IfFileExists{seqsplit.sty}{\usepackage{seqsplit}}{\newcommand{\seqsplit}[1]{#1}}
\protected\def\breakseq#1{\seqsplit{#1}}
\protected\def\breaktt#1{\begingroup\ttfamily\seqsplit{#1}\endgroup}
\setlength{\emergencystretch}{6em}
\tolerance=5000
\hbadness=10000
\hfuzz=1pt
\setlength{\tabcolsep}{2pt}
\AtBeginEnvironment{longtable}{\tiny\renewcommand{\arraystretch}{0.86}\setlength{\tabcolsep}{1pt}}
\AtBeginEnvironment{tabular}{\tiny\renewcommand{\arraystretch}{0.86}\setlength{\tabcolsep}{1pt}}
\AtBeginEnvironment{equation}{\tiny}
\AtBeginEnvironment{equation*}{\tiny}
\AtBeginEnvironment{align}{\tiny}
\AtBeginEnvironment{align*}{\tiny}
\AtBeginEnvironment{itemize}{\footnotesize}
\AtBeginEnvironment{enumerate}{\footnotesize}
\AtBeginEnvironment{description}{\footnotesize}
\setbeamerfont{caption}{size=\tiny}
\setbeamerfont{caption name}{size=\tiny}
\setbeamerfont{normal text}{size=\small}
\setbeamerfont{frametitle}{size=\small}
\setbeamerfont{section title}{size=\footnotesize}
\setbeamerfont{subsection title}{size=\footnotesize}
\setbeamertemplate{section page}{%
  \centering
  \begin{beamercolorbox}[sep=12pt,center,wd=\paperwidth]{section title}
    \parbox{0.86\paperwidth}{\centering\usebeamerfont{section title}\insertsection\par}
  \end{beamercolorbox}
}
\setbeamertemplate{subsection page}{%
  \centering
  \begin{beamercolorbox}[sep=8pt,center,wd=\paperwidth]{subsection title}
    \parbox{0.86\paperwidth}{\centering\usebeamerfont{subsection title}\insertsubsection\par}
  \end{beamercolorbox}
}
\setlength{\abovecaptionskip}{2pt}
\setlength{\belowcaptionskip}{0pt}

% Natbib and cross-reference fallbacks — slides don't load natbib
% or cleveref, but combined-PDF manuscript prose may emit these
% commands. The fallback renders citations as a bracketed key list
% and cross-references as detokenized labels so slides don't fail on
% undefined control sequences or raw underscores. \providecommand is
% a no-op if packages load later.
\providecommand{\citep}[1]{[#1]}
\providecommand{\citet}[1]{#1}
\providecommand{\citealp}[1]{#1}
\providecommand{\citeauthor}[1]{#1}
\providecommand{\citeyear}[1]{#1}
\providecommand{\cref}[1]{\texttt{\detokenize{#1}}}
\providecommand{\Cref}[1]{\texttt{\detokenize{#1}}}
\providecommand{\crefrange}[2]{\texttt{\detokenize{#1}}--\texttt{\detokenize{#2}}}
\providecommand{\Crefrange}[2]{\texttt{\detokenize{#1}}--\texttt{\detokenize{#2}}}
\providecommand{\paragraph}[1]{\textbf{#1}\ }

\newtheorem{proposition}{Proposition}
\newtheorem{hypothesis}{Hypothesis}
\newtheorem{remark}[theorem]{Remark}
\newtheorem{axiom}{Axiom}
\newtheorem{property}{Property}
\usepackage{bookmark}
\IfFileExists{xurl.sty}{\usepackage{xurl}}{} % add URL line breaks if available
\urlstyle{same}
% fallback for those not using the hyperref driver hyperxmp:
\makeatletter
\@ifundefined{xmpquote}{\newcommand{\xmpquote}[1]{#1}}{}
\makeatother
\hypersetup{
  hidelinks,
  pdfcreator={LaTeX via pandoc}}

\author{\texorpdfstring{}{}}
\date{}

\begin{document}

\section{Reproducibility}\label{sec:reproducibility}

\begin{frame}[allowframebreaks]{Reproducibility}
Reproducibility in GNN is not an aspiration layered on top of the
system; it is the operating contract that the pipeline enforces. The
0--24 processing steps are deterministic given a model specification and
a target directory, and every published claim in this manuscript is
traceable to a command that regenerates the underlying artifact. This
section lists only commands that exist in the repository, so that a
reader with a clean checkout can reproduce the pipeline, the validation
gates, and this manuscript itself.
\end{frame}

\begin{frame}[fragile,allowframebreaks]{Pipeline Smoke Run}
\protect\phantomsection\label{pipeline-smoke-run}
The fastest way to confirm a working installation is to drive the full
pipeline over the discrete model family without invoking the optional
LLM steps:

\begin{Shaded}
\begin{Highlighting}[]
\ExtensionTok{uv}\NormalTok{ run python src/gnn/main.py }\AttributeTok{{-}{-}target{-}dir}\NormalTok{ input/gnn\_files/discrete }\AttributeTok{{-}{-}output{-}dir}\NormalTok{ /tmp/gnn{-}smoke }\AttributeTok{{-}{-}skip{-}llm}
\end{Highlighting}
\end{Shaded}

This parses the discrete GNN files, runs visualization and rendering
across the maintained backends, and writes all artifacts under the
chosen output directory. The \texttt{-\/-skip-llm} flag keeps the run
hermetic and free of external API calls: the non-LLM steps all execute,
the steps that would read the skipped LLM outputs record that as a
warning, and the run exits 2 --- the pipeline's documented warning code
(0 success, 1 error, 2 warning) --- rather than 0. To exercise every
\end{frame}

\begin{frame}[fragile,allowframebreaks]{Pipeline Smoke Run}
registered family rather than a single one, drive the manifest through
the model-family acceptance gate given below: pointing
\texttt{-\/-target-dir} at \texttt{input/gnn\_files} covers that tree's
11 corpus directories. All 9 registered family target directories lie
inside that tree, so a single invocation reaches every registered
family.
\end{frame}

\begin{frame}[fragile,allowframebreaks]{Validation Gates}
\protect\phantomsection\label{validation-gates}
GNN's reproducibility guarantees rest on a small set of strict,
deterministic gates that bind the manuscript's quantitative claims to
recomputable ledgers. The model-family acceptance gate runs the
maintained families declared in the manifest and fails on any
regression:

\begin{Shaded}
\begin{Highlighting}[]
\ExtensionTok{uv}\NormalTok{ run python scripts/run\_model\_family\_acceptance.py }\DataTypeTok{\textbackslash{}}
  \AttributeTok{{-}{-}manifest}\NormalTok{ input/model\_family\_manifest.json }\DataTypeTok{\textbackslash{}}
  \AttributeTok{{-}{-}output{-}dir}\NormalTok{ output/model\_family\_acceptance }\AttributeTok{{-}{-}strict}
\end{Highlighting}
\end{Shaded}

\end{frame}

\begin{frame}[fragile,allowframebreaks]{Validation Gates}
The semantic-fidelity gate verifies that a parse \texorpdfstring{\ensuremath{\to}}{->} serialize \texorpdfstring{\ensuremath{\to}}{->} parse
round trip preserves variables, edges, dimensions, parameter shapes,
equations, time semantics, and ontology mappings across the 9 model
families; the cross-framework gate profiles the 7 maintained backends
(PyMDP, RxInfer.jl, JAX, NumPyro, PyTorch, ActiveInference.jl, DisCoPy)
--- refusing any framework outside that set --- and records explicit
compatible and unsupported statuses rather than silently degrading. Both
write their ledgers to an output directory of your choosing:

\end{frame}

\begin{frame}[fragile,allowframebreaks]{Validation Gates}
\begin{Shaded}
\begin{Highlighting}[]
\ExtensionTok{uv}\NormalTok{ run python scripts/run\_semantic\_fidelity\_gate.py }\DataTypeTok{\textbackslash{}}
  \AttributeTok{{-}{-}manifest}\NormalTok{ input/model\_family\_manifest.json }\DataTypeTok{\textbackslash{}}
  \AttributeTok{{-}{-}output{-}dir}\NormalTok{ output/semantic\_fidelity }\AttributeTok{{-}{-}strict}
\ExtensionTok{uv}\NormalTok{ run python scripts/run\_cross\_framework\_reliability.py }\DataTypeTok{\textbackslash{}}
  \AttributeTok{{-}{-}manifest}\NormalTok{ input/model\_family\_manifest.json }\DataTypeTok{\textbackslash{}}
  \AttributeTok{{-}{-}output{-}dir}\NormalTok{ output/cross\_framework }\AttributeTok{{-}{-}strict}
\end{Highlighting}
\end{Shaded}

\end{frame}

\begin{frame}[fragile,allowframebreaks]{Validation Gates}
Under \texttt{-\/-strict}, each gate exits non-zero on the first
mismatch, so these commands double as assertions in an automated
reproduction run. Code-quality reproducibility is enforced separately
through the developer command reference: \texttt{just\ lint} runs the
Ruff linter over \texttt{src} and \texttt{scripts}, and the broader
\texttt{just\ quality} recipe chains formatting, terminology,
documentation, type, and security checks for a full pre-commit gate.
\end{frame}

\begin{frame}[fragile,allowframebreaks]{Manuscript Reproducibility}
\protect\phantomsection\label{manuscript-reproducibility}
This manuscript is itself a reproducible artifact. Every quantitative
value in the prose --- the pipeline step count, the family and backend
counts, the source and test inventories --- is a token rather than a
hard-coded literal, and the deterministic producer regenerates all of
them from the tracked files at the current commit:

\begin{Shaded}
\begin{Highlighting}[]
\ExtensionTok{uv}\NormalTok{ run python scripts/z\_generate\_manuscript\_variables.py}
\end{Highlighting}
\end{Shaded}

\end{frame}

\begin{frame}[fragile,allowframebreaks]{Manuscript Reproducibility}
That command recomputes the \{\{\ldots\}\} tokens, persists them to
\breaktt{output/data/manuscript\_variables.json} for audit, and hydrates
the manuscript sources into \breaktt{output/manuscript/}. The
manuscript's own figures are rebuilt from the same token map:

\begin{Shaded}
\begin{Highlighting}[]
\ExtensionTok{uv}\NormalTok{ run python }\AttributeTok{{-}m}\NormalTok{ scripts.manuscript\_build\_figures}
\end{Highlighting}
\end{Shaded}

The hydrated sources are then rendered to PDF by the docxology
template's render stage. That stage lives in a separate checkout, with
this repository symlinked into it at
\breaktt{projects/active/GeneralizedNotationNotation}; run from the
template root:

\end{frame}

\begin{frame}[fragile,allowframebreaks]{Manuscript Reproducibility}
\begin{Shaded}
\begin{Highlighting}[]
\ExtensionTok{uv}\NormalTok{ run }\AttributeTok{{-}{-}frozen}\NormalTok{ python scripts/pipeline/stage\_03\_render.py }\DataTypeTok{\textbackslash{}}
  \AttributeTok{{-}{-}project}\NormalTok{ GeneralizedNotationNotation}
\end{Highlighting}
\end{Shaded}

The render needs a LaTeX installation providing the packages listed in
\breaktt{manuscript/preamble.md} plus \texttt{seqsplit}; the template
guards \texttt{seqsplit} with \texttt{\textbackslash{}IfFileExists}, so
a missing copy degrades rather than failing the build.

\end{frame}

\begin{frame}[fragile,allowframebreaks]{Manuscript Reproducibility}
Because the variables file is regenerated before rendering, the counts
in the rendered PDF track the repository state at the commit recorded in
\breaktt{output/data/manuscript\_variables.json} (dd195067f): a code
change that alters, for example, the test inventory (373 test files,
4155 test functions) propagates into the prose on the next regeneration
without any manual editing.
\end{frame}

\begin{frame}[fragile,allowframebreaks]{Reproducibility Contract}
\protect\phantomsection\label{reproducibility-contract}
\end{frame}

\begin{frame}[fragile,allowframebreaks]{Reproducibility Contract}
\begin{itemize}
\tightlist
\item
  Do not cite results that cannot be regenerated or directly traced to a
  command in this repository.
\item
  Keep generated outputs under \texttt{output/} and maintained
  manuscript source under \texttt{manuscript/}; treat everything in
  \texttt{output/} as disposable and regeneratable.
\item
  Express every quantitative claim in the prose as a double-brace
  \texttt{\{\{...\}\}} token substituted by
  \breaktt{scripts/z\_generate\_manuscript\_variables.py}, never as a
  hard-coded number.
\item
  Keep private data, credentials, and unpublished sensitive details out
  of the manuscript and out of version control.
\item
  Record the exact verification commands --- the smoke run, the
  acceptance and fidelity gates, and \texttt{just\ lint} --- before
  marking this manuscript publication-ready.
\end{itemize}
\end{frame}

\end{document}
